\chapter{Der Stand im Jahr 2026}\label{kap:aktuell}

\begin{merksatz}
\boxtitel{Was gilt und was behauptet wird}
\begin{center}
{\small
\renewcommand{\arraystretch}{1.5}
\begin{tabular}{@{}P{3.6cm}P{2.0cm}P{2.4cm}@{}}
\toprule
\textbf{Aussage} & \textbf{Beleglage} & \textbf{Urteil}\\
\midrule
Goroutinen sind billig & \textcolor{gruen}{eindeutig} & \textcolor{gruen}{Faktor $\approx4000$}\\
Kleine Schnittstellen tragen & \textcolor{gruen}{gut} & \textcolor{gruen}{Kap.~\ref{kap:schnittstellen}}\\
Gegendruck ohne Puffer & \textcolor{gruen}{gut} & \\
Werkzeuge werden benutzt & \textcolor{gruen}{gut} & \textcolor{gruen}{\textbf{weil mitgeliefert}}\\
Auslieferung ist einfach & \textcolor{gruen}{eindeutig} & \\
\midrule
\enquote{Nebenläufigkeit gelöst} & \textcolor{akzent}{\textbf{falsch}} & \textcolor{akzent}{\textbf{Kap.~\ref{kap:schiefgeht}}}\\
\enquote{Einfach heißt einfach} & \textcolor{akzent}{falsch} & \textcolor{akzent}{Komplexität wandert}\\
\enquote{Fehlerbehandlung ist gelöst} & \textcolor{akzent}{umstritten} & \textbf{$30\,\%$ der Zeilen}\\
\bottomrule
\end{tabular}}
\end{center}

\begin{center}
\begin{minipage}{0.88\textwidth}
\kursiv{Die vierte Zeile ist der praktisch wertvollste Befund dieses Bandes und keine Sprachfrage: Werkzeuge, die mitgeliefert werden und ohne Einrichtung laufen, werden benutzt. Zwei Stunden Einrichtung gegen sechzehn --- das ist die Lehre für jedes Vorhaben, auch in anderen Sprachen.}
\end{minipage}
\end{center}
\end{merksatz}

\begin{aktuell}
\boxtitel{Fünf Entwicklungen}
\textbf{Erstens: Generizität ist angekommen.}

\begin{center}
\renewcommand{\arraystretch}{1.45}
\begin{tabular}{@{}P{3.4cm}P{2.2cm}P{2.2cm}@{}}
\toprule
\textbf{Wirkung} & \textbf{Befund} & \\
\midrule
Bibliotheken & \textcolor{gruen}{\textbf{schlanker}} & \code{slices}, \code{maps}\\
Anwendungscode & \textcolor{mittel}{wenig verändert} & \textcolor{mittel}{selten nötig}\\
Befürchtete Überkomplexität & \textcolor{gruen}{\textbf{blieb aus}} & \\
Übersetzungszeit & \textcolor{mittel}{leicht gestiegen} & \\
\bottomrule
\end{tabular}
\end{center}

\begin{center}
\begin{minipage}{0.86\textwidth}
\kursiv{Die zweite Zeile ist bemerkenswert: Vier Jahre nach der Einführung schreibt der durchschnittliche Anwendungscode kaum eigene Typparameter. Der Nutzen liegt in den Bibliotheken --- genau dort, wo die zwölf Jahre Wartezeit am meisten wehgetan hatten.}
\end{minipage}
\end{center}

\medskip
\textbf{Zweitens: Die Schleifenvariable ist behoben.}

\begin{center}
\renewcommand{\arraystretch}{1.45}
\begin{tabular}{@{}P{3.0cm}P{2.2cm}P{2.6cm}@{}}
\toprule
\textbf{Bis Go 1.21} & \textbf{Seit 1.22} & \\
\midrule
Eine Variable je Schleife & \textbf{eine je Durchlauf} & \\
Häufigster Fehler & \textcolor{gruen}{\textbf{beseitigt}} & \textcolor{gruen}{Kap.~\ref{kap:werkzeuge}}\\
Rückwärts\-verträglich & \textcolor{mittel}{über Version geregelt} & \\
\bottomrule
\end{tabular}
\end{center}

\begin{center}
\kursiv{Eine seltene Entscheidung: Eine Sprachsemantik wurde geändert, weil sie mehr Fehler erzeugte als Nutzen brachte. Die Verträglichkeit wird über die in der Moduldatei genannte Sprachversion geregelt --- alter Code bleibt gültig.}
\end{center}

\medskip
\textbf{Drittens: Regulierung erreicht auch Go.}

\begin{center}
\renewcommand{\arraystretch}{1.45}
\begin{tabular}{@{}P{3.4cm}P{1.8cm}P{2.6cm}@{}}
\toprule
\textbf{Anforderung} & \textbf{Erfüllt?} & \\
\midrule
Speichersicherheit & \textcolor{gruen}{\textbf{ja}} & \textcolor{gruen}{empfohlene Sprache}\\
Stückliste & \textcolor{gruen}{eingebaut} & \textbf{im Binärabbild}\\
Reproduzierbares Bauen & \textcolor{gruen}{\textbf{ja}} & \\
Abhängigkeits\-prüfung & \textcolor{gruen}{ja} & \code{govulncheck}\\
\bottomrule
\end{tabular}
\end{center}

\begin{center}
\begin{minipage}{0.86\textwidth}
\kursiv{Die zweite Zeile ist eine Go-Eigenheit mit regulatorischem Wert: Jede Binärdatei enthält die Liste ihrer Bestandteile samt Versionen und Prüfsummen. Für ein Medizinprodukt bedeutet das, dass die Stückliste nicht gepflegt werden muss --- sie ist im Ergebnis enthalten.}
\end{minipage}
\end{center}

\medskip
\textbf{Viertens: Die Nachfrage bleibt konzentriert.}

\begin{center}
\renewcommand{\arraystretch}{1.45}
\begin{tabular}{@{}P{3.0cm}P{2.2cm}P{2.6cm}@{}}
\toprule
\textbf{Bereich} & \textbf{Anteil} & \\
\midrule
Dienste im Netz & \textcolor{gruen}{\textbf{hoch}} & \\
Werkzeuge, Betrieb & \textcolor{gruen}{hoch} & \\
Anwendungen mit Oberfläche & \textcolor{akzent}{gering} & \\
Wissenschaft & \textcolor{akzent}{\textbf{sehr gering}} & \\
\bottomrule
\end{tabular}
\end{center}

\begin{center}
\kursiv{Fünfzehn Jahre nach der Veröffentlichung ist die Verteilung stabil --- Go hat sich in seiner Nische durchgesetzt und ist darüber hinaus kaum gewachsen. Das ist ein Erfolg für eine Sprache, die für eine Nische entworfen wurde.}
\end{center}

\medskip
\textbf{Fünftens: Werkzeuge, die Go schreiben.}

\begin{center}
\renewcommand{\arraystretch}{1.45}
\begin{tabular}{@{}P{3.0cm}P{2.2cm}P{2.6cm}@{}}
\toprule
\textbf{Wirkung} & \textbf{Befund} & \\
\midrule
Erster Entwurf & \textcolor{gruen}{schneller} & \textbf{wenige Formen}\\
Fehler\-behandlung & \textcolor{gruen}{\textbf{gut erzeugt}} & \textcolor{gruen}{gleichförmig}\\
Nebenläufigkeit & \textcolor{akzent}{\textbf{fehlerhaft}} & \textcolor{akzent}{\textbf{Lecks, fehlender Abbruch}}\\
Prüfbedarf & \textcolor{akzent}{steigt} & \\
\bottomrule
\end{tabular}
\end{center}

\begin{center}
\begin{minipage}{0.86\textwidth}
\kursiv{Die zweite und dritte Zeile stehen im Gegensatz: Die gleichförmige Fehlerbehandlung wird gut erzeugt, weil sie wenige Formen hat. Nebenläufiger Code wird schlecht erzeugt, weil die Fehler nicht im Text sichtbar sind --- eine fehlende Abbruchbehandlung sieht aus wie funktionierender Code.}
\end{minipage}
\end{center}

\begin{center}
\kursiv{Umso wichtiger sind Wettlaufprüfer und Leckprüfung aus Kapitel~\ref{kap:werkzeuge}. Sie prüfen genau das, was beim Lesen nicht auffällt.}
\end{center}
\end{aktuell}

\begin{merksatz}
\boxtitel{Drei Einsichten}
\textbf{Erstens: Weglassen ist eine Entscheidung mit Preis.}

\begin{center}
{\small
\renewcommand{\arraystretch}{1.45}
\begin{tabular}{@{}P{2.8cm}P{2.6cm}P{2.6cm}@{}}
\toprule
\textbf{Auslassung} & \textbf{Gewinn} & \textbf{Preis}\\
\midrule
Ausnahmen & \textbf{sichtbarer Pfad} & \textcolor{akzent}{\textbf{$30\,\%$ der Zeilen}}\\
Vererbung & \textbf{keine Hierarchien} & \textcolor{akzent}{Wiederholung}\\
Überladung & Lesbarkeit & \textcolor{akzent}{Namenswucher}\\
Kleine Sprache & \textbf{$2$--$5$ Tage} & \textcolor{akzent}{wenig Ausdruck}\\
\bottomrule
\end{tabular}}
\end{center}

\begin{center}
\kursiv{Vier Zeilen, acht Spalteneinträge, und keiner davon ist umsonst. Wer nur die mittlere Spalte liest, hält Go für überlegen; wer nur die rechte liest, für armselig. Beide Spalten sind wahr.}
\end{center}

\textbf{Zweitens: Nebenläufigkeit wird ausdrückbar, nicht sicher.}

\begin{center}
\renewcommand{\arraystretch}{1.45}
\begin{tabular}{@{}P{3.0cm}P{2.8cm}P{2.0cm}@{}}
\toprule
\textbf{Fehlerklasse} & \textbf{In Go} & \\
\midrule
Wettlauf & \textcolor{akzent}{möglich} & \textcolor{gruen}{Prüfer findet}\\
Verklemmung & \textcolor{akzent}{möglich} & \textcolor{mittel}{teils erkannt}\\
Goroutinenleck & \textcolor{akzent}{\textbf{leichter als in C}} & \textcolor{akzent}{\textbf{kein Werkzeug}}\\
Speicherfehler & \textcolor{gruen}{\textbf{ausgeschlossen}} & \textcolor{gruen}{der echte Gewinn}\\
\bottomrule
\end{tabular}
\end{center}

\begin{center}
\begin{minipage}{0.88\textwidth}
\kursiv{Die letzte Zeile ist der Fortschritt gegenüber Band 21 --- die häufigste Fehlerklasse überhaupt verschwindet. Die dritte Zeile ist der Rückschritt, den man dafür in Kauf nimmt, und er wird selten genannt.}
\end{minipage}
\end{center}

\textbf{Drittens: Das Handwerk bleibt dasselbe.}

\begin{center}
\renewcommand{\arraystretch}{1.5}
\begin{tabular}{@{}P{3.0cm}P{2.2cm}P{2.6cm}@{}}
\toprule
\textbf{Aus Band 20} & \textbf{Gilt hier?} & \\
\midrule
Aufgabe verstehen & \textcolor{gruen}{\textbf{unverändert}} & \\
Daten zuerst & \textcolor{gruen}{unverändert} & \\
Fehlersuche als Verfahren & \textcolor{gruen}{unverändert} & \textbf{plus Profile}\\
Fremden Code lesen & \textcolor{gruen}{\textbf{leichter}} & \textcolor{gruen}{\textbf{kleine Sprache}}\\
Wann aufhören & \textcolor{gruen}{plus zwei Kriterien} & \\
\bottomrule
\end{tabular}
\end{center}

\begin{center}
\kursiv{Dreimal dasselbe Ergebnis über drei Bände: Fünf von sechs Fertigkeiten sind sprachunabhängig. Nur die vierte Zeile fällt hier besser aus --- eine kleine Sprache ist leichter zu lesen, und das war ihr erklärter Zweck.}
\end{center}

\begin{warnung}
\boxtitelr{Was dieser Band nicht behauptet}
\begin{center}
\renewcommand{\arraystretch}{1.5}
\begin{tabular}{@{}P{3.6cm}P{2.2cm}P{2.0cm}@{}}
\toprule
\textbf{Behauptung} & \textbf{Aufgestellt?} & \\
\midrule
Go ist die beste Sprache & \textcolor{akzent}{\textbf{nein}} & \textbf{Kap.~\ref{kap:wann}}\\
Weglassen ist immer richtig & \textcolor{akzent}{\textbf{nein}} & \textcolor{akzent}{vier Preise}\\
Nebenläufigkeit ist gelöst & \textcolor{akzent}{nein} & \\
Diese Zahlen sind Konstanten & \textcolor{akzent}{nein} & \textbf{Größen\-ordnungen}\\
\bottomrule
\end{tabular}
\end{center}

\begin{center}
\begin{minipage}{0.88\textwidth}
\kursiv{Die zweite Zeile ist die wichtigste. Weglassen ist eine Entscheidung, keine Tugend --- und jede der acht Auslassungen hat einen Anwendungsbereich, außerhalb dessen sie schadet. Kapitel~\ref{kap:auslassungen} nennt für jede das Gegenbeispiel.}
\end{minipage}
\end{center}
\end{warnung}
\end{merksatz}

\bigskip
\begin{flushright}
\itshape\color{grau}
Eine Sprache mit wenigen Mitteln\\
zwingt dazu, das Problem zu verstehen.\\
Eine mit vielen erlaubt es,\\
das Verständnis zu vertagen.\\[10pt]
\normalfont\color{grau}\upshape
--- die Bilanz dieses Bandes.\\
\itshape Beides hat seinen Preis.\\
Nur der erste wird sofort bezahlt.
\end{flushright}
\vspace{6pt}
\noindent Die Bei\-trä\-ge zu Typ\-pa\-ra\-me\-tern\zit{26} erklären die Rück\-nah\-me einer Aus\-las\-sung; Cox\zit{4} do\-ku\-men\-tiert die Ent\-wurfs\-ent\-schei\-dun\-gen; die Spe\-zi\-fi\-ka\-tion\zit{11} bleibt die ver\-bind\-li\-che Quelle.
